
engendrant la filtration canonique.

Ecrivant @@M0@@, avec @@M1@@, et les @@M2@@ étant des monômes par rapport aux @@M3@@, en utilisant l'additivité de @@M4@@, on voit qu'il suffit de faire la démonstration pour @@M5@@, avec @@M6@@ et @@M7@@.

En utilisant le lemme 4.11. et le @@M8@@-@@M9@@-homomorphisme de @@M10@@ dans @@M11@@ qui envoie @@M12@@ sur @@M13@@, on voit qu'il suffit de faire la démonstration pour un élément du type @@M14@@.

Enfin, le polynôme @@M15@@ étant sans composantes isobares de poids @@M16@@, on est ramené au cas @@M17@@, qui résulte lui-même du fait que les polynômes @@M18@@ sont sans composantes isobares de poids @@M19@@ (3.6).

\textbf{Remarques 4.13.} i) La décomposition de @@M20@@ donnée en 4.12 est une décomposition en somme directe de @@M21@@-modules. En d'autres termes, si @@M22@@, et si on écrit @@M23@@, avec @@M24@@, @@M25@@, et @@M26@@, alors @@M27@@.

ii) L'énoncé 4.12. reste vrai pour tout quotient de @@M28@@ par un @@M29@@-idéal contenu dans l'idéal d'augmentation.

iii) L'énoncé analogue à 4.12. obtenu en remplaçant l'indéterminée @@M30@@ par une famille d'indéterminées @@M31@@ est encore vrai.

\textbf{Corollaire 4.14.} @@M32@@.

On définit une flèche de la droite vers la gauche en associant à tout élément @@M33@@, où @@M34@@ est la classe dans @@M35@@ d'un élément @@M36@@ de @@M37@@ et @@M38@@, la classe dans @@M39@@ de @@M40@@. On définit une flèche en sens inverse par décomposition d'après la proposition 4.12 et la remarque 4.13.i). Il est clair que ces deux flèches sont inverses l'une de l'autre.

